\streszczeniepracy{%
  Ta praca doktorska składa się z~serii publikacji, które wprowadzają nowatorskie metody renderowania neuronowego przy użyciu ograniczonych informacji, koncentrując się na Neuronowych Polach Radiancji (NeRF) i~Rozsmarowywaniu Gaussów 3D (3DGS).
  Bada ona, w~jaki sposób te modele konstruują reprezentacje 3D z~obrazów 2D i
  demonstruje sposoby warunkowania tych reprezentacji do generowania wysokiej
  jakości renderów obiektów 3D.

  Przedstawiamy, że zarówno reprezentacje NeRF, jak i~3DGS nie są
  \textit{kontrolowalne} w~interpretowalny sposób---po wytrenowaniu nie można
  już zmienić zawartości sceny tak, aby wynik odzwierciedlał zamierzone zmiany w
  danych wejściowych.
  Proponujemy cztery komplementarne techniki, które wykorzystują proste,
  interpretowalne dane wejściowe pochodzące z~rzadkich danych treningowych i
  rozszerzają te metody, aby skutecznie działały w~scenariuszach uczenia się z
  kilku przykładów.

  Zaczynamy od analizy obszaru Neuronowych Pól Radiancji, zajmując się
  ograniczeniami istniejących podejść i~przedstawiając wkład w~osiągnięcie
  kontrolowalnych pól radiancji.
  Poprzez włączenie częściowych rzadkich danych podczas treningu, nasze
  rozszerzenie NeRF---\conerf---wykorzystuje gładkość sieci neuronowych do
  tworzenia kontrolowalnych obrazów wysokiej jakości dla różnych typów scen.
  Twierdzimy, że rzadkie, ręczne adnotacje, które można uzyskać w~ciągu minuty,
  są wystarczające do zapewnienia interpretowalnej kontroli nad wynikiem NeRFów
  dzięki tendencji sieci neuronowych do faworyzowania niskich częstotliwości.
  Wspieramy ten argument obszernymi wynikami ewaluacji, pokazując, że \conerf
  jest jednym z~pierwszych kroków w~kierunku w~pełni kontrolowalnych pól
  radiancji.

  Aby rozwiązać problem polegania na obszernych, wysokiej jakości adnotacjach
  danych z~filmów wieloperspektywicznych, wprowadzamy nową metodę, \blendfields,
  do trenowania neuronowych pól radiancji w~środowiskach z~wielu kamer z~kilkoma
  przykładami.
  To podejście uczy się wewnętrznych szablonów deformacji, które są płynnie
  mieszane w~czasie wnioskowania, znacznie poprawiając jakość obrazu w
  porównaniu do istniejących punktów odniesienia i~umożliwiając efektywne
  renderowanie z~ograniczonej liczby obrazów wejściowych.
  Wykorzystujemy naszą reprezentację objętościową w~formie siatki czworościanów.
  Używając tradycyjnych technik w~grafice komputerowej, obliczamy wielkości
  deformacji oparte na fizyce, które zapewniają płynną i~interpretowalną
  interpolację między znanymi wyrażeniami.
  Pokazujemy w~naszych eksperymentach, że \blendfields osiąga wyższą jakość
  renderów z~zaledwie kilku ujęć, co kontrastuje z~poprzednimi pracami, które
  opierają się na setkach klatek w~pojedynczym ujęciu.

  Ponadto prezentujemy model do kontrolowania pól radiancji poprzez oświetlenie
  środowiskowe dla ujęć ,,w warunkach naturalnych'' przy różnych warunkach
  oświetleniowych---\lumigauss.
  Poprzez włączenie wstępnie obliczanego transferu promienistego dobrze znanego
  w~grafice komputerowej, model ten umożliwia fizycznie wiarygodne przeliczanie
  oświetlenia sceny i~zapewnia użytkownikom intuicyjną kontrolę nad oświetleniem
  w~zrekonstruowanych scenach.
  Osiągamy ten cel mimo faktu, że każdy obraz treningowy jest wykonywany w
  innych warunkach oświetleniowych i~nie zawiera żadnych informacji o~świetle
  obecnym w~scenie.

  Wiele zastosowań renderowania neuronowego wymaga wysokiej wydajności, dlatego
  w~ostatniej pracy tej rozprawy koncentrujemy się na adaptowalnej efektywności
  obliczeniowej z~naszym \clog.
  Stosuje on strategię uczenia od zgrubnego do dokładnego dla Rozsmarowywania
  Gaussów 3D, która skaluje w~górę siatkę 2D przechowującą reprezentacje ukryte
  Gaussów.
  Ta strategia osiąga konkurencyjne wyniki, jednocześnie umożliwiając wdrożenie
  na różnych urządzeniach obliczeniowych z~minimalną utratą jakości.
  Jak wykazano w~naszych eksperymentach, \clog staje się coraz bardziej wydajny
  podczas wnioskowania z~mniejszą liczbą Gaussów.
  zoptymalizowany dla scen skoncentrowanych na obiektach bez tła.

  Te badania rozwijają najnowsze osiągnięcia w~dziedzinie kontrolowalnych
  neuronowych pól radiancji i~rozszerzają ich zastosowanie do scenariuszy
  uczenia się z~kilku przykładów.
  Te innowacje zwiększają możliwości renderowania neuronowego i~innych obiektów
  w~scenie z~ograniczonych informacji i~otwierają nowe kierunki dla przyszłych
  badań w~tej dziedzinie.
} % koniec streszczenia

\slowakluczowe{Renderowanie Sieciami Neuronowymi, Neuronowe Pola Radiancji, Uczenie z~Kilku Ujęć, Renderowanie Awatarów, Warunkowanie Częściową Informacją, Rozsmarowywanie Gaussów}